% !PRINTED 20
\weberpage{20}
%<<CONT
theiler von $P$ vorausgesetzt ist, und aus der Darstellung (10)
der Gruppe $M$ ergiebt sich, dass $M = Q/Q_1$ ist.

Es ist noch zu zeigen, dass, wenn $M$ Normaltheiler von
$P/Q_1$ ist, auch $Q$ Normaltheiler von $P$ ist. Dies folgt so:

Ist $e$ irgend ein Element in $P$ und $Q_1 b_k$ ein Element in $M$,
so folgt aus der Annahme, dass $M$ Normaltheiler von $P/Q_1$ ist:

\[ Q_1 e^{-1} Q_1 b_k Q_1 e = Q_1 b_h, \]

was wir auch nach der Composition der Theile durch die Formel

\weq{14}{Q_1 e^{-1} Q_1 B Q_1 e = Q_1 B}

ausdrücken können. Da nun $Q_1$ Normaltheiler von $P$ ist, so
kann $Q_1$ mit jedem der anderen Factoren vertauscht werden, so
dass aus (14) folgt:

\[ e^{-1} Q_1 B e = Q_1 B, \]

oder

\weq{15}{e^{-1} Q e = Q,}

wodurch ausgedrückt ist, dass $Q$ Normaltheiler von $P$ ist.

Aus dem Theorem 1. und 2. ergiebt sich das Corollar:

\gesperrt{3. Ein Normaltheiler $Q$ von $P$ ist dann und nur
dann ein grösster Normaltheiler, wenn die complementäre
 Gruppe $P/Q$ einfach ist.}

Wenn $Q$ und $Q'$ zwei Normaltheiler von $P$ sind, so ist auch
das symbolische Product $Q Q'$ ein Normaltheiler.

Denn es ist, weil $Q$ ein Normaltheiler ist

\[ Q Q' = Q' Q, \]

also

\[ Q Q' Q Q' = Q Q Q' Q' = Q Q', \]

und dies ist nach \S.\,4, 1. das Kennzeichen, dass $Q Q'$ eine Gruppe
ist. Dass es ein Normaltheiler von $P$ ist, ergiebt sich dann
ebenfalls nach \S.\,4, 2. aus

\[ Q Q' B = Q B Q' = B Q Q', \]

wenn $B$ irgend ein Theil von $P$ ist.

Ist nun $Q$ ein \gesperrt{grösster Normaltheiler} von $P$, und $Q'$ von
$Q$ verschieden, so ist $Q Q'$ ein Normaltheiler von $P$, der sowohl
$Q$ als $Q'$ in sich enthält und also von $Q$ verschieden ist.
Folglich ist, da es über $Q$ keinen Normaltheiler von $P$ giebt,
ausser $P$ selbst,

\weq{16}{Q Q' = P,}

und ebenso muss auch

\weq{17}{Q' Q = P}
%CONT>>
